% TODO make hw stylesheet
\documentclass[10pt]{scrartcl}
\usepackage[a4paper, total={6in, 8in}]{geometry}
\usepackage[usenames,dvipsnames,svgnames]{xcolor}
\usepackage[shortlabels]{enumitem}
\usepackage[framemethod=TikZ]{mdframed}
\usepackage{amsmath,amssymb,amsthm}
\usepackage[colorlinks]{hyperref}
\usepackage{multicol}
\usepackage{tabularx}

\hypersetup{
	colorlinks=true,
	linkcolor=blue,
	filecolor=magenta,      
	urlcolor=magenta,
	pdftitle={MAT 362 HW}
}

\usepackage{microtype}
\usepackage{mathtools}
\usepackage[headsepline]{scrlayer-scrpage}
\usepackage{thmtools}
\usepackage[textsize=footnotesize]{todonotes}

\mdfdefinestyle{mdbluebox}{roundcorner=10pt,innerbottommargin=9pt,
	linecolor=blue,backgroundcolor=TealBlue!5,}
\declaretheoremstyle[headfont=\sffamily\bfseries\color{MidnightBlue},
mdframed={style=mdbluebox},]{thmbluebox}

\mdfdefinestyle{mdbloobox}{frametitlefont=\bfseries,innerbottommargin=8pt,
	nobreak=true,backgroundcolor=TealBlue!5,linecolor=blue,}
\declaretheoremstyle[headfont=\bfseries\color{MidnightBlue},
mdframed={style=mdbloobox},headpunct={\\[3pt]},postheadspace=0pt,]{thmbloobox}

\mdfdefinestyle{mdredbox}{frametitlefont=\bfseries,innerbottommargin=8pt,
	nobreak=true,backgroundcolor=Salmon!5,linecolor=RawSienna,}
\declaretheoremstyle[headfont=\bfseries\color{RawSienna},
mdframed={style=mdredbox},headpunct={\\[3pt]},postheadspace=0pt,]{thmredbox}

\mdfdefinestyle{mdgreenbox}{linecolor=ForestGreen,backgroundcolor=ForestGreen!5,
	linewidth=2pt,rightline=false,leftline=true,topline=false,bottomline=false,}
\declaretheoremstyle[headfont=\bfseries\sffamily\color{ForestGreen!70!black},
mdframed={style=mdgreenbox},headpunct={ --- },]{thmgreenbox}

\mdfdefinestyle{mdorangebox}{linecolor=RedOrange,backgroundcolor=RedOrange!5,
	linewidth=2pt,rightline=false,leftline=true,topline=false,bottomline=false,}
\declaretheoremstyle[headfont=\bfseries\sffamily\color{RedOrange!70!black},
mdframed={style=mdorangebox},headpunct={ --- },]{thmorangebox}

\mdfdefinestyle{mdblackbox}{linecolor=black,backgroundcolor=RedViolet!5!gray!5,
	linewidth=3pt,rightline=false,leftline=true,topline=false,bottomline=false,}
\declaretheoremstyle[mdframed={style=mdblackbox}]{thmblackbox}

\declaretheoremstyle[spaceabove=6pt,spacebelow=6pt]{basehead}

\declaretheorem[style=basehead,name=Problem]{problem}
\declaretheorem[style=basehead,name=Problem,numbered=no]{problem*}
\declaretheorem[style=thmbloobox,name=Setup]{setup} % for config geo
\declaretheorem[style=thmredbox,name=Example,sibling=problem]{example}
\declaretheorem[style=thmredbox,name=$\bigstar$ Problem,sibling=problem]{niceprob}
\declaretheorem[style=thmbluebox,name=Theorem,sibling=problem]{theorem}
\declaretheorem[style=thmbluebox,name=Corollary,sibling=theorem]{corollary}
\declaretheorem[style=thmbluebox,name=Proposition,sibling=theorem]{prop}
\declaretheorem[style=thmbluebox,name=Lemma,numberwithin=problem]{lemma}
\declaretheorem[style=thmbluebox,name=Theorem,numbered=no]{theorem*}
\declaretheorem[style=thmbluebox,name=Lemma,numbered=no]{lemma*}
\declaretheorem[style=thmgreenbox,name=Claim,numberwithin=problem]{claim}
\declaretheorem[style=thmgreenbox,name=Claim,numbered=no]{claim*}
\declaretheorem[style=thmblackbox,name=Remark,numberwithin=problem]{remark}
\declaretheorem[style=thmblackbox,name=Remark,numbered=no]{remark*}
\declaretheorem[style=thmgreenbox,name=Definition,sibling=theorem]{definition}
\declaretheorem[style=thmgreenbox,name=Definition,numbered=no]{definition*}
\declaretheorem[style=thmorangebox,name=Warning,sibling=theorem]{warning}
\declaretheorem[style=thmorangebox,name=Warning,numbered=no]{warning*}
\declaretheorem[style=thmorangebox,name=Note,numbered=no]{note}
\declaretheorem[style=thmblackbox,name=Exercise,sibling=problem]{exercise}
\declaretheorem[style=thmblackbox,name=Exercise,numbered=no]{exercise*}
\declaretheorem[style=thmblackbox,name=Question,sibling=problem]{question}
\declaretheorem[style=thmblackbox,name=Question,numbered=no]{question*}

\addtolength{\textheight}{3.14cm}
\providecommand{\ol}{\overline}
\providecommand{\ul}{\underline}
\providecommand{\eps}{\varepsilon}
\providecommand{\half}{\frac{1}{2}}
\providecommand{\dang}{\measuredangle} %% Directed angle
\providecommand{\CC}{\mathbb C}
\providecommand{\FF}{\mathbb F}
\providecommand{\NN}{\mathbb N}
\providecommand{\QQ}{\mathbb Q}
\providecommand{\RR}{\mathbb R}
\providecommand{\ZZ}{\mathbb Z}
\providecommand{\dg}{^\circ}
\providecommand{\ii}{\item}
\providecommand{\setdiff}{\smallsetminus}
\providecommand{\alert}[1]{{\sffamily\textbf{\textcolor{blue}{#1}}}}
\providecommand{\inv}{^{-1}}
\providecommand{\mb}[1]{\mathbf{#1}}
\providecommand{\dif}{\;d}
\newcommand*{\horzbar}{\rule[.5ex]{2.5ex}{0.5pt}}

\reversemarginpar

\newenvironment{soln}{\begin{proof}[Solution]}{\end{proof}}
\newenvironment{walkthrough}{\noindent \textbf{Walkthrough.}}{\medskip}
\newlist{walk}{enumerate}{3}
\setlist[walk]{label=\bfseries (\alph*)}

\providecommand{\pow}{\operatorname{Pow}}
\providecommand{\vocab}[1]{{\textbf{\textcolor{ForestGreen}{#1}}}}
\providecommand{\lcm}{\operatorname{lcm}}
\providecommand{\abs}[1]{\left|#1\right|}

\usepackage{asymptote}
\begin{asydef}
	defaultpen(fontsize(10pt));
	unitsize(1cm);
	size(8cm); // set a reasonable default
	usepackage("amsmath");
	usepackage("amssymb");
	settings.tex="pdflatex";
	settings.outformat="pdf";
	// Replacement for olympiad+cse5 which is not standard
	import geometry;
	// recalibrate fill and filldraw for conics
	void filldraw(picture pic = currentpicture, conic g, pen fillpen=defaultpen, pen drawpen=defaultpen)
	{ filldraw(pic, (path) g, fillpen, drawpen); }
	void fill(picture pic = currentpicture, conic g, pen p=defaultpen)
	{ filldraw(pic, (path) g, p); }
	// some geometry
	pair foot(pair P, pair A, pair B) { return foot(triangle(A,B,P).VC); }
	pair orthocenter(pair A, pair B, pair C) { return orthocentercenter(A,B,C); }
	pair centroid(pair A, pair B, pair C) { return (A+B+C)/3; }
	// cse5 abbreviations
	path CP(pair P, pair A) { return circle(P, abs(A-P)); }
	path CR(pair P, real r) { return circle(P, r); }
	pair IP(path p, path q) { return intersectionpoints(p,q)[0]; }
	pair OP(path p, path q) { return intersectionpoints(p,q)[1]; }
	path Line(pair A, pair B, real a=0.6, real b=a) { return (a*(A-B)+A)--(b*(B-A)+B); }
	// cse5 more useful functions
	picture CC() {
		picture p=rotate(0)*currentpicture;
		currentpicture.erase();
		return p;
	}
	pair MP(Label s, pair A, pair B = plain.S, pen p = defaultpen) {
		Label L = s;
		L.s = "$"+s.s+"$";
		label(L, A, B, p);
		return A;
	}
	pair Drawing(Label s = "", pair A, pair B = plain.S, pen p = defaultpen) {
		dot(MP(s, A, B, p), p);
		return A;
	}
	path Drawing(path g, pen p = defaultpen, arrowbar ar = None) {
		draw(g, p, ar);
		return g;
	}
\end{asydef}

\usepackage{mathrsfs}

\ihead{\footnotesize MAT 362 HW10}
\ohead{\footnotesize Last compiled \today}

\title{\vspace{-2em}MAT 362 HW10}
\author{\large Aditya Pahuja}
\date{\large Due 04/24}
\begin{document}
\maketitle

\begin{problem*}4.4.7.
	Intersect the cylinder $x^2+y^2=1$ with a plane passing through
	the $x$-axis and making an angle of $\theta$ with the $xy$-plane.
	\begin{enumerate}[(a)]
		\ii Show that the intersecting curve is an ellipse $C$.
		\ii Compute the absolute value of the geodesic curvature of $C$
		in the cylinder at the points where $C$ meets its axes.
	\end{enumerate}
\end{problem*}

\begin{soln}
	(a) Consider the basis of the plane given by $v_1 = (1,0,0)$
	and the orthogonal vector $v_2 = (0,\cos\theta,\sin\theta)$.
	Given a point $(x,y,z)$ on $C$, we have
	$(x,y,z) = xv_1 + \frac{y}{\cos\theta}v_2$
	(since $y/\cos\theta$ is the only choice of coefficient
	that matches with the second component of $(x,y,z)$).
	We have that $x^2 + y^2 =1$ a priori, so, within the intersecting plane,
	$C$ has equation $(x')^2 + (y')^2\cos^2\theta = 1$,
	where points $(x',y')$ are respectively the $v_1$, $v_2$-coordinates
	with respect to the orthonormal basis $\{v_1,v_2\}$ of the plane;
	this is exactly the equation of an ellipse.
	
	(b) At a given point $p$, since
	$k^2 = k_g^2 + k_n^2$, we have that $|k_g| = k\sin\phi$
	where $\phi\in[0,\pi)$ is the angle between the normal to the curve
	and normal to the surface at $p$.
	
	The axes of $C$ intersect $C$ at two points on the $x$-axis,
	where $\phi = 0$ (hence the geodesic curvature is $0$),
	and two points on the $yz$-plane, where $\phi = \theta$.
	In the latter case we wish to compute the curvature of the ellipse at
	$(0,1,\tan\theta)$ and $(0,-1,-\tan\theta)$; since both values
	are identical by symmetry we just compute for the former.
	We can parametrize the ellipse as $r(t) = (\cos t,\sin t/\cos\theta,0)$
	(with respect to orthonormal basis $v_1$, $v_2$, $v_1\times v_2$)
	so that the curvature at $t=\frac\pi2$ is
	\[\frac{|r'\times r''|}{|r'|^3} = \frac{|(0,0,1/\cos\theta)|}{|(-1,0,0)^3|}
	= \sec\theta.\]
	The geodesic curvature is thus $\tan\theta$.
\end{soln}

\begin{problem*}4.4.11.
	State precisely and prove: the algebraic value of the
	covariant derivative is invariant under orientation-preserving isometries.
\end{problem*}

\begin{soln}
	Statement:
	Let $w$ be a differentiable field of unit vectors along a parametrized curve
	$\alpha\colon I\to S$ on an oriented surface $S$.
	Then, if $f$ is an isometry of $S$, the algebraic value of
	$f\circ w$ (which is a vector field on $f\circ\alpha$)
	is equal to the algebraic value of $w$.
	We thus want to show that if
	\[\frac{Dw}{dt} = \lambda(N\wedge w(t)),\quad
	\frac{D(f\circ w)}{dt} = \lambda'(f(N)\wedge f(w(t))\]
	then $\lambda = \lambda'$.
	Being an isometry, $\langle N,w(t)\rangle = \langle f(N),f(w(t))\rangle$,
	as well as $|N| = |f(N)|$ and $|w(t)| = |f(w(t))|$, so
	$|N\wedge w(t)| = |f(N)\wedge f(w(t))$; thus it suffices to show that
	the magnitudes of $Dw/dt$ and $D(f\circ w)/dt$ are the same.
	
	On the other hand $Dw/dt$ is just the projection of $w'$
	on $T_{\alpha(0)}(S)$ and $D(f\circ w)/dt$ is the projection of
	$(f\circ w)'$ on $T_{f(\alpha(0))}(S)$,
	but $f$ being an isometry means the latter is the same length as the former
	(``the geometric configurations of the objects are the same'' roughly),
	which gives us the result.
\end{soln}

\begin{problem*}4.4.18.
	Consider a geodesic which starts at a point $p$ in the upper part
	of a hyperboloid of revolution $x^2+y^2-z^2=1$
	and makes an angle $\theta$ with the parallel passing through $p$
	in such a way that $\cos\theta = 1/r$, where $r$ is
	the distance from $p$ to the $z$-axis.
	Show that by following the geodesic in the direction of
	decreasing parallels, it approaches asymptotically
	the parallel $x^2+y^2=1$, $z=0$.
\end{problem*}

\begin{soln}
	Parametrize the geodesic by arc length starting at $p$
	to get the map $\alpha\colon[0,\infty)\to H$ where $H$ is the hyperboloid.
	Then let $r(s)$ be the distance to the $z$ axis
	and $\theta(s)$ the angle between $\alpha'(s)$ and
	the parallel passing through $\alpha(s)$. Then,
	\[r(s)\cos \theta(s) = r(0)\cos\theta(0) = r\cdot1/r = 1\]
	for all $s$.
	Now let $c = \inf_{s\ge 0}r(s) = \lim_{s\to\infty} r(s)$.
	We know that $\sqrt{1-\frac1c} < \sin\theta(s)$ for all $s$
	and, if $\alpha(s) = (x(s),y(s),z(s))$, then $z'(s) = \sin\theta(s)$
	because $|\alpha'(s)| = 1$.
	Thus if $c > 1$, given $\eps > 0$, pick $t_\eps$ so that
	$r(t_\eps) < c+\eps$. Because of the bound on $z'$,
	$z(t+1) \le z(t)-\eps$. Clearly $z > 0$ for all $s$
	(otherwise we cross through the $xy$-plane, when $r(s)=1<c$),
	but this is a contradiction, since we can just take $n$ large enough
	so that $z(t+n) \le z(t) - n\eps < 0$.
	
	That is, $c>1$ leads to a contradiction,
	so the limiting value of $r$ must be $1$ (since of course $r(s)\ge1$)
	and the limiting value of $\cos\theta(s)$ must also be $1$
	because of Clairaut's relation, which is what we want.
\end{soln}

\begin{problem*}4.4.23.
	Show that the isometries of the unit sphere $S^2$ are the restrictions
	to $S^2$ of the linear orthogonal transformations of $\RR^3$.
\end{problem*}

\begin{soln}
	Let $f$ be an isometry of $S^2$.
	Then $f$ maps great circles to great circles
	(since it maps geodesics to geodesics,
	and the great circles are all the geodesics of $S^2$)
	and preserves the ``angular distance'' between two points on $S^2$
	(i.e. $\langle p,q\rangle = \langle f(p),f(q)\rangle$).
	Extending this map to $f(tp) = tf(p)$ for $p\in S^2$ and real $t$,
	we can write, for arbitrary $t_1p_1,t_2p_2\in\RR^3$
	(with $p_1,p_2\in S^2$), that $\langle t_1p_2,t_2p_2\rangle =
	\langle t_1f(p_1),t_2f(p_2)\rangle$, i.e. $f$ preserves
	inner products (because distances and angles are preserved),
	which is enough to conclude that it is
	an orthogonal linear transformation by a previous homework problem. 
\end{soln}


















\end{document}